\documentclass[a4paper,12pt]{article}
\usepackage[T2A]{fontenc}
\usepackage[utf8]{inputenc}

\usepackage[russian]{babel}

\usepackage[top=25mm,bottom=30mm,left=20mm,right=20mm]{geometry}

% Красная строка для первого предложения в разделе
\usepackage{indentfirst}

% Равные промежутки между словами и предложениями
\frenchspacing

\usepackage{amsmath,amssymb,mathtools,amsthm}
\usepackage{icomma}

% Стиль нумерации формул
\renewcommand{\theequation}{\thesection.\arabic{equation}}

% Отображать номера только тех формул, на которые есть ссылка
\mathtoolsset{showonlyrefs=true}

% Более привычные греческие буквы
\renewcommand{\epsilon}{\varepsilon}
\renewcommand{\phi}{\varphi}
\renewcommand{\kappa}{\varkappa}

% Знаки сравнения с наклонёнными чёрточками
\renewcommand{\le}{\leqslant}
\renewcommand{\leq}{\leqslant}
\renewcommand{\ge}{\geqslant}
\renewcommand{\geq}{\geqslant}

\renewcommand{\emptyset}{\varnothing}

\newcommand{\ds}{\displaystyle}

% Обозначения числовых множеств
\newcommand{\N}{\mathbb{N}}
\newcommand{\Z}{\mathbb{Z}}
\newcommand{\Q}{\mathbb{Q}}
\newcommand{\R}{\mathbb{R}}

\DeclareMathOperator{\im}{Im}
\DeclareMathOperator{\re}{Re}

% Стрелочки в рамочках для доказательств разных частей двусторонних утверждений
\newcommand{\Boxedrightarrow}{\raisebox{0.3ex}{\(\boxed{\Rightarrow}\)}\;}
\newcommand{\Boxedleftarrow}{\raisebox{0.3ex}{\(\boxed{\Leftarrow}\)}\;}

\newtheoremstyle{main}
    {4pt}                                        % расстояние до
    {4pt}                                        % расстояние после
    {}                                           % шрифт тела
    {\parindent}                                 % размер отступа красной строки
    {\bfseries}                                  % шрифт шапки
    {.}                                          % пунктуация после заголовка
    { }                                          % пунктуация после шапки
    {\thmname{#1}\thmnumber{ #2}\thmnote{ (#3)}} % параметры шапки
\theoremstyle{main}

\newtheorem{definition}{Определение}[section]

\newtheorem{theorem}{Теорема}[section]

% Квадратики вокруг доказательства теоремы
\AtBeginDocument{\renewcommand*{\proofname}{\(\square\)\nopunct}}
\renewcommand{\qedsymbol}{\(\blacksquare\)}

\newtheorem*{corollary}{Следствие}
\newtheorem*{statement}{Утверждение}
\newtheorem*{addition}{Дополнение}
\newtheorem*{note}{Замечание}
\newtheorem*{anote}{Замечание автора}
\newtheorem*{notation}{Обозначение}
\newtheorem*{reminder}{Напоминание}
\newtheorem*{example}{Пример}
\newtheorem*{examples}{Примеры}

% Команда для переноса во встраиваемых в строку формулах. Взята из книги
% «Набор и вёрстка в системе LaTeX» С.М. Львовского.
\newcommand*{\hm}[1]{#1\nobreak\discretionary{}%
    {\hbox{$\mathsurround=0pt #1$}}{}}

\usepackage{enumitem}

% Стиль списка для математических свойств
\newenvironment{properties}[1][]
    {\begin{enumerate}[label=\arabic*\(^{\circ}\).,#1]}
    {\end{enumerate}}

\sloppy

% Стили ссылок
\usepackage{hyperref}
\usepackage[usenames,dvipsnames,svgnames,table,rgb]{xcolor}
\hypersetup{
    unicode=true,
    colorlinks=true,
    linkcolor=black!15!blue,
    citecolor=green,
    filecolor=NavyBlue,
    urlcolor=NavyBlue,
}

% Настройка колонтикулов
\usepackage{titleps}

\newpagestyle{toc}{
    \setheadrule{0.4pt}
    \sethead{Многомерный анализ, интегралы и ряды}{}{}
    \setfoot{}{\thepage}{}
}

\newpagestyle{main}{
    \setheadrule{0.4pt}
    \sethead{Многомерный анализ, интегралы и ряды}{}
        {\hyperlink{toc}{\;Назад к содержанию}}
    \setfoot{}{\thepage}{}
}

\pagestyle{main}

\usepackage{tikz}
\usetikzlibrary{decorations.pathreplacing}

\newcommand\Section[1]{
    \refstepcounter{section}
    \section*{\raggedright
        \Roman{section}. #1}
    \addcontentsline{toc}{section}{
        \Roman{section}. #1}
}

\newcommand\Subsection[1]{
    \refstepcounter{subsection}
    \subsection*{\raggedright
        \S \arabic{subsection}. #1}
    \addcontentsline{toc}{subsection}{
        \S \arabic{subsection}. #1}
}

\usepackage{float}